\begin{tomtat}
	\begin{minipage}[t]{0.3\textwidth}
		\vspace{0pt}
		\begin{tomtat}
			\begin{center}
				\textbf{ĐỀ 12}
			\end{center}
			
		\end{tomtat}
	\end{minipage}
	\hfill
	\begin{minipage}[t]{0.5\textwidth}
		\vspace{0pt}
		\begin{tomtat}
			\begin{center}
				\includegraphics[scale=0.08]{images/12}
			\end{center}
		\end{tomtat}
	\end{minipage}
\end{tomtat}
\textbf{I. PHẦN TRẮC NGHIỆM}
\begin{ex}
	Phương trình $(x-2)(x-3)=0$ có nghiệm là:
	\choice{$x \in \{2;-3\}$}{$x \in \{-2;3\}$}{$x \in \{-2;-3\}$}{$x \in \{2;3\}$}
	\loigiai{
		Ta có phương trình tích $(x-2)(x-3)=0$.
		Một tích bằng 0 khi và chỉ khi ít nhất một trong các thừa số của nó bằng 0.
		Do đó, ta có hai trường hợp:
		\begin{itemize}
			\item $x-2=0 \Rightarrow x=2$
			\item $x-3=0 \Rightarrow x=3$
		\end{itemize}
		Vậy, tập nghiệm của phương trình là $S = \{2; 3\}$.
		\\ \textbf{ĐÁP ÁN D}
	}
\end{ex}

\begin{ex}
	Điều kiện xác định của biểu thức $\sqrt{2x-4}$ là:
	\choice{$x \le 4$}{$x \ge 4$}{$x \ge 2$}{$x \le 2$}
	\loigiai{
		Biểu thức $\sqrt{2x-4}$ xác định khi biểu thức dưới dấu căn không âm.
		Tức là $2x-4 \ge 0$.
		\begin{align*}
			2x-4 &\ge 0 \\
			2x &\ge 4 \\
			x &\ge \frac{4}{2} \\
			x &\ge 2
		\end{align*}
		Vậy điều kiện xác định của biểu thức là $x \ge 2$.
		\\ \textbf{ĐÁP ÁN C}
	}
\end{ex}

\begin{ex}
	Giá trị biểu thức $\frac{1}{2-\sqrt{3}}+\frac{1}{2+\sqrt{3}}$ là:
	\choice{4}{3}{2}{1}
	\loigiai{
		Ta quy đồng mẫu số:
		\begin{align*}
			A &= \frac{1(2+\sqrt{3})}{(2-\sqrt{3})(2+\sqrt{3})} + \frac{1(2-\sqrt{3})}{(2+\sqrt{3})(2-\sqrt{3})} \\
			&= \frac{2+\sqrt{3}}{2^2 - (\sqrt{3})^2} + \frac{2-\sqrt{3}}{2^2 - (\sqrt{3})^2} \\
			&= \frac{2+\sqrt{3}}{4 - 3} + \frac{2-\sqrt{3}}{4 - 3} \\
			&= \frac{2+\sqrt{3}}{1} + \frac{2-\sqrt{3}}{1} \\
			&= 2+\sqrt{3} + 2-\sqrt{3} \\
			&= 4
		\end{align*}
		\\ \textbf{ĐÁP ÁN A}
	}
\end{ex}

\begin{ex}
	Với giá trị nào của m thì đường thẳng $(d) y=(m^2-4)x+2m-3$ và đường thẳng $y=5x+3$ song song với nhau?
	\choice{$m=3$}{$m=-3$}{$m=\pm3$}{$m=\sqrt{3}$}
	\loigiai{
		Hai đường thẳng $y=ax+b$ và $y=a'x+b'$ song song với nhau khi và chỉ khi $a=a'$ và $b \ne b'$.
		Đường thẳng $(d)$ có $a = m^2-4$ và $b = 2m-3$.
		Đường thẳng $y=5x+3$ có $a'=5$ và $b'=3$.
		Để hai đường thẳng song song, ta cần:
		1. $a=a' \Rightarrow m^2-4 = 5 \Rightarrow m^2 = 9 \Rightarrow m = \pm 3$.
		2. $b \ne b' \Rightarrow 2m-3 \ne 3$.
		
		Xét trường hợp $m=3$:
		$b = 2(3)-3 = 6-3 = 3$.
		Vì $b=b' (3=3)$, nên khi $m=3$, hai đường thẳng trùng nhau, không song song.
		
		Xét trường hợp $m=-3$:
		$b = 2(-3)-3 = -6-3 = -9$.
		Vì $b \ne b' (-9 \ne 3)$, nên khi $m=-3$, hai đường thẳng song song.
		
		Vậy giá trị của $m$ để hai đường thẳng song song là $m=-3$.
		\\ \textbf{ĐÁP ÁN B}
	}
\end{ex}

\begin{ex}
	Trong mặt phẳng tọa độ Oxy cho parabol $(P) y=x^2$ và đường thẳng $(d) y=2x-m+3$. Giá trị m để đường thẳng $(d)$ cắt $(P)$ tại 2 điểm phân biệt là:
	\choice{$m \le 4$}{$m \ge 4$}{$m < 4$}{$m > 4$}
	\loigiai{
		Phương trình hoành độ giao điểm của $(P)$ và $(d)$ là:
		$x^2 = 2x-m+3$
		$x^2 - 2x + m - 3 = 0 \quad (*)$
		Để đường thẳng $(d)$ cắt $(P)$ tại 2 điểm phân biệt thì phương trình $(*)$ phải có hai nghiệm phân biệt.
		Điều này xảy ra khi biệt thức $\Delta' > 0$.
		Phương trình $(*)$ có $a=1, b'=-1, c=m-3$.
		$\Delta' = (b')^2 - ac = (-1)^2 - 1(m-3) = 1 - m + 3 = 4 - m$.
		Để có hai nghiệm phân biệt, $\Delta' > 0$:
		$4 - m > 0$
		$4 > m$
		$m < 4$.
		\\ \textbf{ĐÁP ÁN C}
	}
\end{ex}

\begin{ex}
	Cho tam giác ABC vuông tại A có AB=3 và AC=4. Khi đó độ dài đoạn thẳng BC bằng:
	\choice{1}{25}{7}{5}
	\loigiai{
		Trong tam giác ABC vuông tại A, theo định lý Pitago, ta có:
		$BC^2 = AB^2 + AC^2$
		$BC^2 = 3^2 + 4^2$
		$BC^2 = 9 + 16$
		$BC^2 = 25$
		$BC = \sqrt{25} = 5$.
		(Vì độ dài đoạn thẳng là số dương)
		\\ \textbf{ĐÁP ÁN D}
	}
\end{ex}
\begin{ex} % Câu 7
	Tính thể tích $V$ của hình cầu có bán kính $R=3\text{cm}$.
	\choice{$V=180\pi \text{ cm}^3$}{$V=9\pi \text{ cm}^3$}{$V=72\pi \text{ cm}^3$}{$V=36\pi \text{ cm}^3$}
	\loigiai{
		Thể tích của hình cầu được tính bằng công thức: $V = \frac{4}{3}\pi R^3$.
		Với $R=3 \text{ cm}$:
		\begin{align*}
			V &= \frac{4}{3}\pi (3)^3 \\
			&= \frac{4}{3}\pi \cdot 27 \\
			&= 4\pi \cdot \frac{27}{3} \\
			&= 4\pi \cdot 9 \\
			&= 36\pi (\text{cm}^3)
		\end{align*}
		\\ \textbf{ĐÁP ÁN D}
	}
\end{ex}

\begin{ex} % Câu 8
	Bạn Nam gieo một con xúc xắc 10 lần liên tiếp thì thấy mặt 4 chấm xuất hiện 3 lần. Xác suất thực nghiệm xuất hiện mặt 4 chấm là:
	\choice{$\frac{4}{10}$}{$\frac{3}{10}$}{$\frac{7}{10}$}{$\frac{3}{14}$}
	\loigiai{
		Xác suất thực nghiệm của một biến cố được tính bằng tỉ số giữa số lần biến cố đó xảy ra và tổng số lần thực hiện phép thử.
		Số lần gieo xúc xắc (tổng số lần thực hiện phép thử) là $N = 10$ lần.
		Số lần mặt 4 chấm xuất hiện (số lần biến cố xảy ra) là $k = 3$ lần.
		Xác suất thực nghiệm xuất hiện mặt 4 chấm là:
		$$ P(\text{mặt 4 chấm}) = \frac{k}{N} = \frac{3}{10} $$
		\\ \textbf{ĐÁP ÁN B}
	}
\end{ex}
\newpage
\textbf{II. PHẦN TỰ LUẬN}
\begin{ex} % Câu 9
	Rút gọn biểu thức: $B = \left(\frac{1}{\sqrt{x}} + \frac{1}{3-\sqrt{x}}\right) : \frac{1}{3-\sqrt{x}}$ với $x > 0, x \ne 9$.
	\loigiai{
		Với $x > 0, x \ne 9$, ta có:
		\begin{align*}
			B &= \left(\frac{3-\sqrt{x} + \sqrt{x}}{\sqrt{x}(3-\sqrt{x})}\right) : \frac{1}{3-\sqrt{x}} \\
			&= \frac{3}{\sqrt{x}(3-\sqrt{x})} \cdot \frac{3-\sqrt{x}}{1}
		\end{align*}
		Vì $x \ne 9$ nên $3-\sqrt{x} \ne 0$. Ta có thể rút gọn $3-\sqrt{x}$.
		\begin{align*}
			B &= \frac{3}{\sqrt{x}}
		\end{align*}
		Vậy $B = \frac{3}{\sqrt{x}}$.
	}
\end{ex}

\begin{ex} % Câu 10 (Hệ phương trình)
	Giải hệ phương trình sau: $\begin{cases} 2x-3y=7 \\ x+5y=-3 \end{cases}$.
	\loigiai{
		Ta có hệ phương trình:
		$$ \begin{cases} 2x-3y=7 & (1) \\ x+5y=-3 & (2) \end{cases} $$
		Từ phương trình (2), suy ra $x = -3 - 5y$.
		Thay $x = -3 - 5y$ vào phương trình (1):
		\begin{align*}
			2(-3-5y) - 3y &= 7 \\
			-6 - 10y - 3y &= 7 \\
			-13y &= 7 + 6 \\
			-13y &= 13 \\
			y &= -1
		\end{align*}
		Thay $y=-1$ vào $x = -3 - 5y$:
		\begin{align*}
			x &= -3 - 5(-1) \\
			x &= -3 + 5 \\
			x &= 2
		\end{align*}
		Vậy nghiệm của hệ phương trình là $(x; y) = (2; -1)$.
	}
\end{ex}

\begin{ex} % Câu 10 (Phương trình bậc hai)
	(1,5 điểm). Cho phương trình $x^2+3x-2=0$. Tính $T = \frac{x_1^2}{x_2} + \frac{x_2^2}{x_1} - 5(x_2+x_1)$.
	\loigiai{
		Phương trình đã cho là $x^2+3x-2=0$.
		Đây là phương trình bậc hai có dạng $ax^2+bx+c=0$ với $a=1, b=3, c=-2$.
		Ta tính biệt thức $\Delta = b^2-4ac = (3)^2 - 4(1)(-2) = 9 + 8 = 17$.
		Vì $\Delta = 17 > 0$ nên phương trình có hai nghiệm phân biệt $x_1, x_2$.
		Vì $c=-2 \ne 0$, nên $x_1, x_2 \ne 0$.
		
		Theo hệ thức Vi-ét, ta có:
		$$ \begin{cases} 
			x_1+x_2 = \dfrac{-b}{a} = \dfrac{-3}{1} = -3 \\ 
			x_1x_2 = \dfrac{c}{a} = \dfrac{-2}{1} = -2 
		\end{cases} $$
		Biến đổi biểu thức $T$:
		\begin{align*}
			T &= \frac{x_1^3 + x_2^3}{x_1x_2} - 5(x_1+x_2) \\
			&= \frac{(x_1+x_2)(x_1^2 - x_1x_2 + x_2^2)}{x_1x_2} - 5(x_1+x_2) \\
			&= \frac{(x_1+x_2)((x_1+x_2)^2 - 2x_1x_2 - x_1x_2)}{x_1x_2} - 5(x_1+x_2) \\
			&= \frac{(x_1+x_2)((x_1+x_2)^2 - 3x_1x_2)}{x_1x_2} - 5(x_1+x_2)
		\end{align*}
		Thay giá trị $x_1+x_2 = -3$ và $x_1x_2 = -2$ vào biểu thức $T$:
		\begin{align*}
			T &= \frac{(-3)((-3)^2 - 3(-2))}{-2} - 5(-3) \\
			&= \frac{-3(9 + 6)}{-2} + 15 \\
			&= \frac{-3 \cdot 15}{-2} + 15 \\
			&= \frac{-45}{-2} + 15 \\
			&= \frac{45}{2} + \frac{30}{2} \\
			&= \frac{75}{2}
		\end{align*}
		Vậy giá trị của biểu thức $T$ là $\frac{75}{2}$.
	}
\end{ex}

\begin{ex} % Câu 11
	Một hộp đựng 4 thẻ, mỗi thẻ được ghi một trong các số 1, 2, 3, 4; hai thẻ khác nhau thì ghi hai số khác nhau. Rút ngẫu nhiên lần lượt 2 thẻ từ hộp đó (mỗi lần rút thẻ ra không bỏ thẻ lại hộp). Tính xác suất của biến cố A: "Rút được 2 thẻ có tổng hai số ghi trên thẻ là một số chia hết cho 3".
	\loigiai{
		Tập hợp tất cả các kết quả có thể xảy ra khi rút lần lượt 2 thẻ từ 4 thẻ (không hoàn lại) là các cặp số có thứ tự:
		$\Omega = \{$
		(1,2), (1,3), (1,4), 
		(2,1), (2,3), (2,4), 
		(3,1), (3,2), (3,4), 
		(4,1), (4,2), (4,3) 
		$\}$.
		Số phần tử của không gian mẫu là $n(\Omega) = 12$.\\
		
		Tập hợp các kết quả thuận lợi cho biến cố A: "Rút được 2 thẻ có tổng hai số ghi trên thẻ là một số chia hết cho 3".
		Các cặp có tổng chia hết cho 3 là:
		\begin{itemize}
			\item Tổng bằng 3: (1,2), (2,1)
			\item Tổng bằng 6: (2,4), (4,2)
		\end{itemize}
		Vậy, tập hợp các kết quả thuận lợi cho biến cố A là: $A = \{(1,2), (2,1), (2,4), (4,2)\}$.
		Số phần tử của biến cố A là $n(A) = 4$.
		
		Xác suất của biến cố A là:
		$$ P(A) = \frac{n(A)}{n(\Omega)} = \frac{4}{12} = \frac{1}{3} $$
	}
\end{ex}

\begin{ex} % Câu 12
	(1,5 điểm): Cho một thửa ruộng hình chữ nhật, biết rằng nếu chiều rộng tăng thêm 2m, chiều dài giảm đi 2m thì diện tích thửa ruộng đó tăng thêm 30m$^2$; và nếu chiều rộng giảm đi 2m, chiều dài tăng thêm 5m thì diện tích thửa ruộng giảm đi 20m$^2$. Biết rằng một mùa vụ 1m$^2$ trên thửa ruộng đó thu hoạch được 0,8kg thóc. Tính khối lượng thóc thu hoạch được trong một mùa vụ của thửa ruộng lúc ban đầu.
	\loigiai{
		Gọi chiều dài của thửa ruộng lúc ban đầu là $x$ (m) và chiều rộng là $y$ (m).
		Điều kiện: $x > 0, y > 0$.
		Diện tích ban đầu của thửa ruộng là $S_1 = xy$ (m$^2$).
		
		Trường hợp 1: Chiều rộng tăng thêm 2m, chiều dài giảm đi 2m thì diện tích tăng thêm 30m$^2$.
		Chiều rộng mới: $y+2$ (m).
		Chiều dài mới: $x-2$ (m). (Điều kiện $x>2$)
		Diện tích mới: $S' = (x-2)(y+2)$.
		Ta có phương trình: $S' - S_1 = 30$
		\begin{align*}
			(x-2)(y+2) - xy &= 30 \\
			xy + 2x - 2y - 4 - xy &= 30 \\
			2x - 2y - 4 &= 30 \\
			2x - 2y &= 34 \\
			x - y &= 17 \quad (1)
		\end{align*}
		
		Trường hợp 2: Chiều rộng giảm đi 2m, chiều dài tăng thêm 5m thì diện tích giảm đi 20m$^2$.
		Chiều rộng mới: $y-2$ (m). (Điều kiện $y>2$)
		Chiều dài mới: $x+5$ (m).
		Diện tích mới: $S'' = (x+5)(y-2)$.
		Ta có phương trình: $S_1 - S'' = 20$
		\begin{align*}
			xy - (x+5)(y-2) &= 20 \\
			xy - (xy - 2x + 5y - 10) &= 20 \\
			xy - xy + 2x - 5y + 10 &= 20 \\
			2x - 5y &= 10 \quad (2)
		\end{align*}
		
		Từ (1) và (2) ta có hệ phương trình:
		$$ \begin{cases} x - y = 17 \\ 2x - 5y = 10 \end{cases} $$
		Từ (1) $\Rightarrow x = y+17$. Thay vào (2):
		\begin{align*}
			2(y+17) - 5y &= 10 \\
			2y + 34 - 5y &= 10 \\
			-3y &= 10 - 34 \\
			-3y &= -24 \\
			y &= 8
		\end{align*}
		Thay $y=8$ vào $x = y+17$:
		$x = 8+17 = 25$.
		Kiểm tra điều kiện: $x=25>2$, $y=8>2$. Thỏa mãn.
		
		Diện tích ban đầu của thửa ruộng là $S_1 = xy = 25 \cdot 8 = 200$ (m$^2$).
		Khối lượng thóc thu hoạch được là: $200 \text{ m}^2 \times 0,8 \text{ kg/m}^2 = 160$ kg.
		
		Vậy khối lượng thóc thu hoạch được trong một mùa vụ của thửa ruộng lúc ban đầu là 160 kg.
	}
\end{ex}
\begin{ex} 
	Cho $\triangle ABC$ có ba góc nhọn nội tiếp đường tròn $(O;R)$. Các đường cao $AD, BF, CE$ của $\triangle ABC$ cắt nhau tại $H$.
	\begin{enumerate}
		\item[a.] Chứng minh tứ giác $BEHD$ nội tiếp một đường tròn.
		\item[b.] Kéo dài $AD$ cắt đường tròn $(O)$ tại điểm thứ hai $K$. Kéo dài $KE$ cắt đường tròn $(O)$ tại điểm thứ hai $I$. Kẻ $OM$ vuông góc với $BC$ tại $M$. Gọi $P$ là tâm đường tròn ngoại tiếp $\triangle AEF$. Gọi $N$ là giao điểm của $CI$ và $EF$. Chứng minh $CE^2 = CN \cdot CI$ và ba điểm $M, N, P$ thẳng hàng.
	\end{enumerate}
	\loigiai{
		\begin{center}
			\begin{tikzpicture}[line join = round, line cap=round,>=stealth,font=\footnotesize,scale=1.5]
				\usetikzlibrary{calc,intersections,angles}
				
				\def\r{3}
				
				\coordinate (O) at (0,0);
				\coordinate (A) at (110:\r);
				\coordinate (B) at (-160:\r);
				\coordinate (C) at (-20:\r);
				
				\coordinate (E) at ($(A)!(C)!(B)$);
				\coordinate (F) at ($(A)!(B)!(C)$);
				\coordinate (D) at ($(B)!(A)!(C)$);
				\coordinate (G) at ($(A)!(E)!(C)$);
				
				\path[name path=lineCE_H] (C) -- ($(C)!1.5!(E)$);
				\path[name path=lineBF_H] (B) -- ($(B)!1.5!(F)$);
				\path[name intersections={of=lineCE_H and lineBF_H, by=H}];
				
				\coordinate (K) at ($2*(D)-(H)$);
				
				\coordinate (P) at ($(A)!0.5!(H)$);
				\coordinate (M) at ($(B)!0.5!(C)$);
				\coordinate (N) at ($(E)!0.5!(F)$);
				
				\path[name path=lineCN_I] (C) -- ($(C)!2.5!(N)$);
				\path[name path=lineKE_I] (K) -- ($(K)!2.5!(E)$);
				\path[name intersections={of=lineCN_I and lineKE_I, by=I}];
				
				\draw (O) circle (\r);
				\draw (A) -- (B) -- (C) -- cycle;
				\draw (B) -- (F);
				\draw (C) -- (E)--(G);
				\draw (A) -- (K) -- (I)--(C);
				\draw (O) -- (M);
				\draw (E) -- (F);
				\draw (P) -- (M);
				\draw (A) -- (D);
				
				\fill (O) circle (1pt);
				\fill (A) circle (1pt);
				\fill (B) circle (1pt);
				\fill (C) circle (1pt);
				\fill (E) circle (1pt);
				\fill (F) circle (1pt);
				\fill (D) circle (1pt);
				\fill (G) circle (1pt);
				\fill (H) circle (1pt);
				\fill (K) circle (1pt);
				\fill (P) circle (1pt);
				\fill (M) circle (1pt);
				\fill (N) circle (1pt);
				\fill (I) circle (1pt);
				
				\node[below left] at (O) {$O$};
				\node[above] at (A) {$A$};
				\node[below left] at (B) {$B$};
				\node[below right] at (C) {$C$};
				\node[above left] at (E) {$E$};
				\node[above right] at (F) {$F$};
				\node[below right] at (D) {$D$};
				\node[right] at (G) {$G$};
				\node[left] at (H) {$H$};
				\node[below] at (K) {$K$};
				\node[right] at (P) {$P$};
				\node[below] at (M) {$M$};
				\node[below] at (N) {$N$};
				\node[left] at (I) {$I$};
				
				\pic[draw, angle radius=3mm]{right angle = A--D--B};
				\pic[draw, angle radius=3mm]{right angle = O--M--C};
				\pic[draw, angle radius=3mm]{right angle = A--E--C};
				\pic[draw, angle radius=3mm]{right angle = B--F--C};
				\pic[draw, angle radius=3mm]{right angle = A--G--E};
				
			\end{tikzpicture}
			
			
		\end{center}
		Chứng minh được tứ giác $AEHF$ nội tiếp
		
		Xét đường tròn ngoại tiếp tứ giác $AEHF$ có
		
		$\widehat{FEH} = \widehat{FAH}$ (hai góc nội tiếp cùng chắn cung $FH$) hay
		
		$\widehat{CEN} = \widehat{KAC}$ (1)
		
		Xét $(O)$ có $\widehat{KAC} = \widehat{KIC}$ (hai góc nội tiếp cùng chắn cung $KC$) hay $\widehat{KAC} = \widehat{EIC}$ (2)
		
		Từ (1) và (2) suy ra $\widehat{CEN} = \widehat{EIC}$
		
		Xét $\triangle CEN$ và $\triangle CIE$ có: $\widehat{ECI}$ : chung; $\widehat{CEN} = \widehat{EIC}$ (cmt)
		
		Nên $\triangle CEN \sim \triangle CIE$ (g-g) Suy ra $\frac{CE}{CI} = \frac{CN}{CE} \Rightarrow CE^2 = CN \cdot CI$ (đpcm)
		
		+ Xét tam giác $OBC$ cân tại $O$
		
		Vì $OM \perp BC$ tại $M$ nên $OM$ là đường cao của tam giác cân nên $OM$ cũng là đường trung tuyến, do đó $M$ là trung điểm $BC$.
		
		Xét $\triangle EBC$ vuông tại $E$ có $M$ là trung điểm $BC$ nên $ME = \frac{1}{2}BC$.
		
		Tương tự ta có $MF = \frac{1}{2}BC$. Do đó $ME = MF (= \frac{1}{2}BC)$ suy ra $M$ thuộc trung trực của $EF$.
		
		Vì $P$ là tâm đường tròn ngoại tiếp $\triangle AEF$ nên $PE = PF$.
		
		Suy ra $P$ thuộc trung trực của $EF$. Vì vậy $PM$ là trung trực của $EF$ (3).
		
		Để chứng minh $M, N, P$ thẳng hàng ta đi chứng minh $N \in PM$.
		
		Kẻ $EG \perp AC, G \in AC$.
		
		Áp dụng hệ thức lượng vào $\triangle AEC$ vuông tại $E$, đường cao $EG$, ta có $CE^2 = CG \cdot CA$.
		
		Theo phần b có $CE^2 = CN \cdot CI$ nên $CG \cdot CA = CN \cdot CI \Rightarrow \frac{CG}{CN} = \frac{CI}{CA}$.
		Xét $\triangle CNG$ và $\triangle CAI$ có \\
		$\frac{CG}{CN} = \frac{CI}{CA}$ (cmt) và $\widehat{ICA}$ : chung
		
		Nên $\triangle CNG \sim \triangle CAI$ (c.g.c) $\Rightarrow \widehat{NGC} = \widehat{CIA}$ (hai góc tương ứng) hay $\widehat{NGF} = \widehat{CIA}$ (*)
		
		Xét $(O)$ có $\widehat{CIA} = \widehat{CBA}$ (hai góc nội tiếp cùng chắn $\stackrel{\frown}{CA}$) (**)
		
		Chứng minh được tứ giác $BEFC$ nội tiếp một đường tròn $\Rightarrow \widehat{EBC} + \widehat{EFC} = 180^\circ$ (hai góc đối nhau)
		
		Mà $\widehat{AFE} + \widehat{EFC} = 180^\circ$ (hai góc kề bù) $\Rightarrow \widehat{EBC} = \widehat{AFE}$ hay $\widehat{ABC} = \widehat{NFG}$ (***)
		
		Từ (*), (**), (***) ta suy ra $\widehat{NGF} = \widehat{NFG}$
		
		Do đó $\triangle NGF$ cân tại $N$ suy ra $NG = NF$
		
		Xét $\triangle EGF$ vuông tại $G$ có $\widehat{NGF} = \widehat{NFG}$ nên $\widehat{NGE} = \widehat{NEG}$.
		
		Do đó $\triangle NGE$ cân tại $N$ suy ra $NG = NE$
		
		Khi đó $NE = NF$ hay $N$ là trung điểm $EF$ (4)
		
		Từ (3) và (4) suy ra $N \in PM$ hay ba điểm $M,N,P$ thẳng hàng.
		}
		
		
\end{ex}
